\chapter{On Lying as the Origin of Meaning}
\label{ch:lie}
% ===================================================================

The last chapter asked what a symbol system would have to be and answered with a
functor that organizes the conditions attractively and then fails. This chapter
takes up what the failure leaves standing.

Chapter~\ref{ch:trick} established a negative result and gave it two-sided
evidence. Rongorongo and Linear~A are corpora without populations: enough text to
train a model to emit well-formed strings, and no community of speakers left.
Sperm whale codas are the mirror case --- no Rosetta stone, no cognates, and
real progress, because the population is alive and its behavior can be watched.
The conclusion was that the correspondence between a symbol system and the world
is carried in the community of speakers rather than in the symbol system, and
that decipherment is either transport into an anchor or direct observation of the
using population, neither of which is distributional.

That chapter also insisted, in Observation~\ref{obs:soluble}, that the reference
bootstrap problem is not open in the sense that its solubility is in doubt. It
has been solved twice, independently. And it isolated three features common to
both solutions: the solver was a population, the interface was not chosen, and
there was an external criterion supplying the fitness signal.

What it did not supply was a mechanism. It said that the crossing has been made
and did not say by what motion.

This chapter records a candidate. It is due to Michael Benedikt, it is that
language originates in lying, and i think it is the most interesting thing
anybody has said to me about the bootstrap problem. i want to be clear at the
outset that what follows is speculative in a way that Chapter~\ref{ch:notation}
was not. The propositions are targets. The reason to state them anyway is that
two of them, if they went through, would upgrade the negative result of
Chapter~\ref{ch:trick} from an observation to a theorem, and would say something
sharp and unwelcome about what grounding requires.

\section{Benedikt's argument}
\label{sec:benedikt}

Michael Benedikt was ACSA Distinguished Professor of Architecture at the
University of Texas at Austin, held the Hal Box Chair in Urbanism, and wrote
\emph{Deconstructing the Kimbell} among much else \cite{benedikt1991}. He died in
August 2025. The argument below is recalled from conversation and should be read
as personal communication; a search has surfaced no published version of it.

\paragraph{The setting.} A band of early hominids on the plain, not far from the
trees, dividing food --- fruit, or the spoils of a kill.

\paragraph{The prior state.} Many animals make reflexive noises grounded in
function: the sounds of breathing, howling in pain, the alarm call. These are
triggered by conditions. They are not meaningful in the way that language is
meaningful, and the whole argument turns on what the prior state lacks.

\paragraph{The event.} Eve notices movement in the grass. The conditions land on
her nervous system in just such a way as to trigger the snake warning call. She
calls; the band scatters for the trees. In that split second Eve determines that
it was not a snake --- it was the wind --- and at the same moment notices that
she now stands in an advantageous position with respect to the division of the
food.

What she sees is \emph{the gap}: the distance between the reflexively triggered
sound and the behavior of the band. And she sees that the gap admits
exploitation.

\paragraph{The follow-on.} If Eve does this regularly, others may notice the
correlation between the snake call and the division of the food. The band learns
Eve's trick from Eve's behavior. Language is a virus \cite{burroughs1962}.

\paragraph{The reading.} It is the gap that gives the opening for meaning. A
reflexive call and its triggering condition are welded together, and nothing
there needs interpreting. Prise them apart and semantics becomes both possible
and necessary. This is the sense in which lying is the knowledge of good and evil
offered at the center of the garden --- a remark i will leave as a remark, since
the formal apparatus does not reach it.

\section{The empirical situation is better than anecdotal}
\label{sec:deception-empirical}

The parable has data behind it, and the data sharpen it in a direction the
parable does not anticipate.

\paragraph{Tufted capuchins.} Wheeler reports that wild tufted capuchins give
predator calls in non-predatory contexts where the caller stands to gain
\cite{wheelerbc2009}. False alarms are given by subordinates more than by
dominants, rise with the contestability of the food and with its scarcity, and
are given from positions in which the caller would gain if the others reacted.
Listeners leave the patch; the caller takes the food. This is Eve's scenario,
observed.

\paragraph{Counterdeception.} Listeners are more likely to ignore the same call
type when it is produced during competitive feeding than in other contexts, and
the difference is carried primarily by a drop in escape responses rather than in
vigilance \cite{wheelerbc2010,wheelerbc2013}. The receivers have become
context-sensitive about a signal that used to require no context.

\paragraph{Fork-tailed drongos.} Drongos steal food using \emph{mimicked} alarm
calls of other species, and vary which call they mimic
\cite{flower2011,flower2014}. The reason given is that a deceptive signal loses
power as its frequency rises relative to honest signals, so varying the signal is
what sustains the deception.

\paragraph{Earlier work.} Munn's kleptoparasitic shrike-tanagers and antshrikes
\cite{munn1986}, whose status as alarm rather than aggressive signals has since
been questioned.

\begin{observation}[Deception grows the vocabulary]
\label{lie:obs:vocab}
Under frequency-dependent discounting, deception does not merely exploit an
existing repertoire. It exerts pressure to expand the repertoire, because a
fresh signal starts at the honest response rate.
\end{observation}

Observation~\ref{lie:obs:vocab} is a mechanism for vocabulary growth and not
only for the birth of meaning, and it is the part of the picture the parable
alone would not have produced.

\section{Three stages, and a re-typing}
\label{sec:stages}

The translation into the framework is short, and its shortness is the point:
nothing is added to the calculus.

\paragraph{Stage 0: the reflex.} The call is a rewrite rule, not a symbol. It is
a sensor-driven emission, and the receiver's response is another rewrite. The
correspondence to the world is a fact about the wiring, and neither party needs a
hypothesis. In the terms of Table~1 of \chSci both sides are non-modal: this is
perception, reading structure the world has already computed into legible form.

\paragraph{Stage 1: the decoupling.} Eve's innovation is to make the call channel
writable by her policy independently of the sensor. In the channel classification
of \chSci --- channels enabling internal reductions, channels enabling
interaction with the environment, and channels doing both --- this is a
\emph{re-typing}. An emission driven only from the sensor side becomes driven
from the policy side as well. The gap is a change in which redexes can fire on a
name, and that is all it is.

\paragraph{Stage 2: semantics as the receiver's forced expense.} Once emission is
decoupled, \emph{call heard} no longer entails \emph{snake present}. The
receiver's predicate cannot remain non-modal. It must be indexed by context, and
the context is a hypothesis about the \emph{sender} rather than about the world.
This is machinery the Turn already built: context-labeled modalities, and an
environment whose most interesting contents are other learners.

\begin{observation}[The thesis]
\label{lie:obs:thesis}
Meaning is not a property of the sign. Meaning is the theory that receivers are
forced to build once senders are capable of lying. Semantics is the receiver's
bill.
\end{observation}

\section{Goodman, read backwards}
\label{sec:goodman-backwards}

Chapter~\ref{ch:notation} adopted Goodman's conditions provisionally and spent
them: on the two bisimulations, on what a budget buys, and on the construction
that fails. The adoption was a bet that the conditions were the right ones rather
than merely the first ones available. With Stage~0 in hand there is something to
say about that bet, and it is not what i expected to be saying.

Goodman's \emph{marks} are, in his own gloss, physical inscriptions or
\emph{utterances}. A triggered alarm call is a mark. The apparatus therefore
applies to Stage~0 directly, with no reinterpretation and no analogy, and we can
simply run the five conditions on a reflexive population and see what happens.

\begin{table}[t]
\centering\small
\fitwidth{%
\begin{tabular}{@{}llll@{}}
\toprule
\textbf{Condition} & \textbf{Stage 0 (reflex)} & \textbf{After the lie} &
\textbf{Repaired by} \\
\midrule
Syntactic disjointness (\ref{cond:sd})
  & holds & broken by mimicry & --- \\
Syntactic finite diff.\ (\ref{cond:sfd})
  & holds if discrete & unchanged & Conjecture~\ref{conj:fd} \\
\textbf{Unambiguity} (\ref{cond:ua})
  & \textbf{holds maximally} & \textbf{broken by the lie} & --- \\
Semantic disjointness (\ref{cond:semd})
  & holds & broken by mimicry & --- \\
Semantic finite diff.\ (\ref{cond:semfd})
  & holds if discrete & unchanged & Conjecture~\ref{conj:fd} \\
\bottomrule
\end{tabular}}
\caption{Goodman's five conditions evaluated on a reflexive signalling
population, and again after emission is decoupled from the sensor. The two
conditions the budget move of \S\ref{sec:budget} repairs are exactly the two
that deception leaves alone.}
\label{lie:tab:goodman}
\end{table}

\subsection{A perfect notation that means nothing}
\label{sec:perfect-notation}

Take the conditions in turn on a Stage~0 population with a discrete repertoire.
No mark belongs to two characters, so Condition~\ref{cond:sd} holds. The call
types have no structurally intermediate forms, so Condition~\ref{cond:sfd}
holds. Distinct call types fire on distinct triggering conditions, so
Condition~\ref{cond:semd} holds. And Condition~\ref{cond:ua} --- that each
character have a single compliance class, not different compliants at different
times or in different contexts --- holds \emph{maximally}, because the call and
its condition are welded. That weld is what Benedikt's argument says the prior
state consists of.

\begin{observation}[Notationality is maximized where reference is absent]
\label{lie:obs:perfect}
A discrete reflexive signalling system satisfies all five of
Conditions~\ref{cond:sd}--\ref{cond:semfd}. It is, by
Chapter~\ref{ch:notation}'s adopted criterion, a notational system. It is also
the system in which, on the argument of this chapter, nothing means anything.
\end{observation}

Observation~\ref{lie:obs:perfect} sharpens Remark~\ref{rem:invariance} rather
than contradicting it. That remark said Goodman's conditions cannot
\emph{select} a compliance relation, being invariant under relabeling. The
reflex case says something less comfortable: the conditions are best satisfied
precisely where no compliance relation was selected at all, because the world
wired it. Notational adequacy is not merely insufficient for reference. It is
maximized at the point where reference is absent, and degrades as reference
arrives.

And the condition that degrades is nameable. Stage~1 introduces a second
production path for a mark that had one, so that \emph{call heard} has one
compliant when a snake is present and another when a subordinate stands over
contested food. That is Condition~\ref{cond:ua}, failing, in the condition's own
words. The gap is not analogous to a failure of unambiguity and is not a
consequence of one. It is the same fact under two descriptions.

\subsection{Compliance requires possible non-compliance}
\label{sec:possible-noncompliance}

There is a stronger reading available, and it is the one i would defend.

Goodman's compliance relation presupposes that compliance can fail. The whole
wrong-note paradox of \S\ref{sec:failures} exists because a performance can fail
to be a performance of the work; if it could not, there would be no paradox and
no need for the budget repair. But a trigger cannot fail. If the condition holds
and the wiring is intact the call fires, and if the call fired the condition
held. There is no room for a mis-transcription.

So at Stage~0 there is not a compliance relation that happens to be perfectly
satisfied. There is no compliance relation --- only causation, described twice.
Character and compliant are not two things standing in a relation; they are one
thing under two names, which is why the relation between them cannot fail and
why Condition~\ref{cond:ua} is satisfied for free.

This shows up structurally as an arrow count, and the arrow count is the cleanest
version of the point.

\begin{observation}[The lie makes the roundtrip a roundtrip]
\label{lie:obs:roundtrip}
Goodman's adequacy, by \S\ref{sec:roundtrip}, is a roundtrip: character to
compliant and back, compliant to character and back. At Stage~0 there is one
arrow, world to mark, and no return, because the mark is caused by the compliant
rather than produced from a character. Stage~1 supplies a second producer of
marks, so the mark ceases to determine its producer. Stage~2 requires the
receiver to invert, and by Conjecture~\ref{lie:conj:mixture} the inversion is of
a mixture rather than of a function.
\end{observation}

Read that way, Goodman's framework requires exactly the structure that deception
introduces. Before the lie there is a transduction and no notation; the
transcription direction does not exist to be checked. This is worth holding
alongside \S\ref{sec:transformer} of the last chapter, which located the
transducer inside the calculus and reassigned the network to proposal. The
transducer is not where notation lives, and Stage~0 is a population that has a
transducer and nothing else.

\subsection{Deception attacks what the budget cannot repair}
\label{sec:complementary}

The right-hand columns of Table~\ref{lie:tab:goodman} are worth reading
together, because the division is exact.

Conjecture~\ref{conj:fd} makes the two finite-differentiation conditions
budget-relative, and its reading was that a poorer agent has a coarser notation.
It says nothing about unambiguity or about either disjointness condition, and
this is not an oversight: those three are not statements about how finely an
observer can separate things, so no amount of budget bears on them. A rich agent
and a poor agent face the same ambiguous sign.

Deception breaks exactly those three. The lie breaks unambiguity. Mimicry breaks
the two disjointness conditions --- and does so in a way worth stating plainly,
because it is not a semantic phenomenon at that point. A fork-tailed drongo
emitting a pied babbler's alarm call \cite{flower2011,flower2014} has put one
mark into two characters. Condition~\ref{cond:sd} is syntactic, and deception at
the community level is a syntactic attack on it.

Between the budget move and the deception mechanism, then, all five conditions
are accounted for, and neither account touches the other's. That is not an
argument that either is correct. It is an argument that they are two halves of
one thing, and that Goodman's conditions decompose along a seam nobody had
reason to look for until there was a liar in the picture.

\subsection{What the ethologists have been measuring}
\label{sec:functional-reference}

One check is available immediately, and it is favorable in an unusual way: the
relevant measurements have already been made by people who had never heard of
Goodman.

The standing criteria for \emph{functional reference} in animal signalling are
two \cite{macedonia1993}. Production specificity: the signal is given to one
class of stimulus. And the perception criterion: the signal elicits its response
in the absence of supporting contextual cues --- the property that made the
vervet alarm calls the textbook case \cite{seyfarth1980}. The second is Condition~\ref{cond:ua}, tested from
the receiver's side and stated in behavioral terms. And the live dispute in that
field is whether receivers integrate context after all --- which is precisely
what the counterdeception results of \cite{wheelerbc2010,wheelerbc2013} report
they do, and precisely what Stage~2 predicts they must.

A second hit lands on the syntactic side. Quantitative reanalysis of vervet
alarm calls finds the leopard, eagle and snake calls well separated, with
overlap appearing once calls given in \emph{aggressive} contexts are included
\cite{price2015}. Condition~\ref{cond:sd} fails exactly where the socially
strategic use enters, which is where this chapter says it should.

Neither observation was made to test anything here, which is the reason to
report them.

\subsection{The objection, which i cannot dispose of}
\label{sec:deflation}

A deflationary reading is available and should be stated at full strength.

One may say: at Stage~0 the compliance class of a character simply \emph{is} its
triggering condition class. That is a perfectly good compliance relation. The
system is therefore notational, Observation~\ref{lie:obs:perfect} is correct but
unsurprising, and all it establishes is that notationality is not sufficient for
meaning --- which Remark~\ref{rem:invariance} already said, by a shorter route.
On this reading nothing above is new.

The reply is Observation~\ref{lie:obs:roundtrip}: a theory of notation whose
central failure mode --- the wrong note --- is unstatable in its best-behaved
instance is telling us something about what it measures. But that reply depends
on the claim that compliance requires possible non-compliance, and that claim is
at present an intuition with a paragraph attached rather than an argument. It is
the first thing i would want written down properly, and until it is, the
deflationary reading is live and the stronger one is a proposal.

\subsection{A prediction the pairing generates}
\label{sec:digitality}

Combining the two halves of Table~\ref{lie:tab:goodman} yields a claim that
neither makes alone.

\begin{conjecture}[Deception is a force for digitality]
\label{lie:conj:digital}
Under the frequency-dependent discounting of
Observation~\ref{lie:obs:vocab}, a fresh signal is worth manufacturing only if
it starts at the honest response rate, and it can do that only if receivers do
not confuse it with the discounted ones. The pressure of
Conjecture~\ref{lie:conj:elab} is therefore not merely toward more characters
but toward characters satisfying Conditions~\ref{cond:sd}
and~\ref{cond:sfd}. Deception pressure should drive a repertoire from the
analogue end of Goodman's distinction toward the digital end.
\end{conjecture}

Conjecture~\ref{lie:conj:digital} is comparative and testable. Signalling
systems are documented as graded, as discrete, and as mixed, and the variation is
measurable. The prediction is that discreteness should covary with deception
pressure, and the drongo is a suggestive extreme: it does not manufacture new
characters at all but imports discrete ones wholesale from other species, which
is the cheapest possible way to buy distinguishability.

If it holds, the digital / analogue distinction acquires a second economic
reading to set beside the one Conjecture~\ref{conj:fd} gives it. There, a
notation is as fine as its observer can afford. Here, a notation is as discrete
as its liars have made it necessary to be.

\section{Statements to attempt}
\label{sec:lie-props}

The retrospect of \S\ref{sec:goodman-backwards} named the condition that fails
and left the failure informal. What follows are the targets that would make it
precise, together with the ones that would give the mechanism empirical content.
None is established.

\begin{conjecture}[No semantics before the lie]
\label{lie:conj:before}
In a population where every emission is sensor-driven, the receiver's optimal
policy is a non-modal predicate on the signal: the modal rungs buy nothing.
\end{conjecture}

The corollary would be that an alarm call in a wholly honest population is not a
symbol --- and the framework could say so with a purchase criterion rather than
with a definition, which is the difference between a result and a stipulation.

\begin{conjecture}[The lie forces the modal rung]
\label{lie:conj:forces}
With a nonzero deception rate the non-modal predicate is strictly dominated, and
the receiver's best available response requires context-indexing. The transition
has a threshold in the deception rate.
\end{conjecture}

The threshold should be computable in the currency of the ladder returns of
\chGame, and Conjectures~\ref{lie:conj:before} and~\ref{lie:conj:forces} together
are the pair such a computation could actually decide.

\begin{conjecture}[The deception equilibrium]
\label{lie:conj:equil}
The liar's yield is proportional to the response rate, which declines in the
population deception rate; hence an interior equilibrium.
\end{conjecture}

The empirical predictions are already made and already met: subordinates lie
more, lying rises with contestability, listeners discount contextually
\cite{wheelerbc2009,wheelerbc2010}.

\begin{conjecture}[Elaboration under discounting]
\label{lie:conj:elab}
At equilibrium there is pressure to diversify the signal. Vocabulary size grows
as a function of the deception pressure and of the cost of manufacturing a new
distinguishable name.
\end{conjecture}

Freshness by quotation says names are cheap to manufacture, which predicts fast
growth --- plausibly too fast. What limits it is a real question and a good one.

\begin{conjecture}[The trick is horizontally transmitted code]
\label{lie:conj:horiz}
The band learns Eve's trick by watching Eve, not by descending from her: a quoted
process copied between contemporaries. This is the reproduction material of
\chSci with the vertical arrow replaced by a horizontal one.
\end{conjecture}

Burroughs' virus is not a metaphor under Conjecture~\ref{lie:conj:horiz}. It is
an $R_{0}$, and the thing it spreads degrades the resource it spreads on.

\begin{conjecture}[Strategic use makes the corpus unidentifiable]
\label{lie:conj:mixture}
If emission is policy-driven, the joint distribution over signals and
world-states is a mixture over sender policies whose weights are facts about the
population. A corpus gives the marginal, and the marginal does not identify the
mixture.
\end{conjecture}

Conjecture~\ref{lie:conj:mixture} is the statement with reach outside the
framework, because it says that no corpus of any size recovers the
correspondence, and says it as a claim about strategic signaling rather than as a
complaint about data scarcity. If it goes through, the central observation of
Chapter~\ref{ch:trick} is upgraded: the correspondence lives in the community
\emph{because} the senders are strategic. Linear~A resists because its semantics
was never in the distribution. Whale song yields because the senders remain
available to be modeled.

\begin{conjecture}[Grounding is a drive, not a read]
\label{lie:conj:drive}
A transformer over a corpus can only read. Under
Conjecture~\ref{lie:conj:mixture} the readable object is the wrong one, and the
identifying information is available only to an agent that can emit into the
community and observe the response --- that is, to one that can perform the
intervention separating the mixture components.
\end{conjecture}

This sharpens Section~\ref{sec:transformer} of the last chapter rather than
contradicting it. There the network was reassigned from transducer to proposal
mechanism, on the grounds that the transducer is already exact and a network
cannot improve on it. Conjecture~\ref{lie:conj:drive} says something stronger
about the other half: not merely that a corpus is the wrong place to look for the
transducer, but that a corpus is the wrong \emph{kind} of object for recovering
correspondence at all, for a reason that has nothing to do with its size. If it
holds, grounding requires participation rather than perception, which is a sharp
and unwelcome claim about what embodiment has to mean. Its interest is that it is
stated in the framework's own pricing vocabulary rather than imported.

\section{Caveats, and what the argument does not license}
\label{sec:lie-caveats}

\paragraph{Deception is not the only bootstrap.} Lewis conventions and Skyrms'
signaling games produce meaning by reinforcement under aligned interests, with no
liar anywhere \cite{lewis1969,skyrms2010}. The defensible claim is narrower than
``lying is necessary for signaling''. It is that \emph{lying is what makes
meaning require modeling}, forcing the receiver up from correlation to theory.
The two routes should yield different kinds of symbol system, and characterizing
the difference is a good problem.

\paragraph{Deception is parasitic on an honest baseline.} The reflexive layer
must exist first, or there is nothing to exploit. This constrains any engineered
bootstrap: build a functionally grounded reflex layer before there is anything
for a policy to detach from.

\paragraph{The ethological alignment should be checked rather than reinvented.}
Dawkins and Krebs argue that signals are manipulation rather than information
transfer \cite{dawkinskrebs1978}, which is Benedikt's thesis arriving from
ethology; and the equilibrium of Conjecture~\ref{lie:conj:equil} has presumably
been worked out in the costly-signaling literature \cite{msharper2003} in a
different currency. Import it and pay it in tokens.

\paragraph{Intentionality is not claimed and should not be.} The capuchin
literature is careful about whether false alarms are deliberate, to the point of
testing whether adrenocortical stress explains the calling. The framework does
not need deliberateness. It needs only that emission be policy-driven rather than
sensor-driven, which is a structural condition on channels and is exactly what
the conjectures above are stated in terms of.

\section{Open questions}
\label{sec:lie-open}

\begin{question}
Does the deception rate play the role that the environment plays in the
inversion of \chComp --- so that whether \emph{meaning} is affordable is a
property of the environment rather than of the learner, exactly as whether
\emph{truth} is affordable turned out to be?
\end{question}

\begin{question}
Can Conjecture~\ref{lie:conj:mixture} be given as a genuine identifiability
theorem, with the mixture stated over a space of sender policies and the failure
of identifiability exhibited rather than argued?
\end{question}

\begin{question}
Of the three common features of Observation~\ref{obs:soluble}, the third --- an
external criterion --- is here instantiated as the food share, and the first ---
that the solver was a population --- is here \emph{explained} rather than
observed, since a lie requires someone to lie to. Does the second, that the
interface was not chosen, receive a similar treatment?
\end{question}

\begin{question}
Deception should be impossible in the limit where agents reason over one
another's source rather than over one another's claims. In the three grades of
access of Table~1 of \chSci that limit is the reflection grade, and deception
lives at grades one and two. Does it return under bounded verification, when the
cost of checking a property exceeds the stake --- making obfuscation an economic
attack rather than an informational one? And if two agents are extensionally
indistinguishable on the fragment each can afford to check, deception appears to
live entirely in the complement, which suggests a framework-native definition:
\emph{deception is engineering the opponent into the undecided region and
harvesting their default}.
\end{question}

\section{What is claimed}
\label{sec:lie-claimed}

Claimed: that the reflex-to-symbol transition can be stated in the framework as a
re-typing of a channel, with nothing added to the calculus; that on that reading
meaning is the receiver's forced expense rather than a property of the sign
(Observation~\ref{lie:obs:thesis}); and that frequency-dependent discounting
gives vocabulary growth a mechanism (Observation~\ref{lie:obs:vocab}).

Not claimed: any of the seven conjectures, and in particular not
Conjecture~\ref{lie:conj:mixture}, which is the one i would most like to be
true and therefore the one to be most careful about. It is stated as a target
because the shape of the required proof is visible --- a mixture over sender
policies, and a demonstration that the marginal does not identify the weights ---
and because both sides of it are things somebody could actually work on.

There is a reason this chapter sits where it does. Chapter~\ref{ch:notation}
reached a construction that fails, and named the failure precisely: the calculus
supplies no proposal mechanism, no way to get from a stream of world to a set of
builtins. The temptation at that point is to conclude that what is missing is
more data or a better architecture. The suggestion here is that what is missing
is a liar --- that the information which fixes reference is not in any record of
what was said, because what was said was said by somebody with an interest, and
the only way to learn what an interested speaker means is to speak back and find
out what it costs them.
